\documentclass[10pt,twocolumn,letterpaper]{article}

\usepackage[review]{cvpr}
\usepackage{newtxtext,newtxmath}
\usepackage{graphicx}
\usepackage{amsmath}
\usepackage{amssymb}
\usepackage{booktabs}
\usepackage{enumitem}
\usepackage{xcolor}
\definecolor{cvprblue}{rgb}{0.21,0.49,0.74}
\usepackage[pagebackref,breaklinks,colorlinks,allcolors=cvprblue]{hyperref}

\def\paperID{*****}
\def\confName{CVPR}
\def\confYear{2027}

\begin{document}

\title{Pixel Art Is a Palette, Not a Downscale --- Supplementary Material}
\author{Anonymous CVPR submission\\Paper ID \paperID}
\maketitle

\section{What the body left out}

The body reports what a reader needs to use the method. This document carries the rest: the data pipeline in enough
detail to rebuild the corpus, the results that did not earn a place in the recipe, and the measurements behind claims
the body states in one sentence. Section numbers here are independent of the body.

\section{Corpus construction}
\label{sup:data}

\paragraph{From sheets to sprites.}
Source material arrives as sprite sheets and loose PNGs. Whole images whose longest side is at most 64\,px are kept as
they are --- these are the natively tiny sprites the low-resolution buckets are starving for. Larger images are treated
as sheets and cut into alpha-connected components between 8 and 64\,px. Every cut records the archive, the member file
and the licence it came with.

\paragraph{Filters.}
A candidate is discarded unless its shorter side is at least 10\,px, its aspect ratio is at most 3:1, the fraction of
pixels with $\alpha > 16$ lies in $[0.15, 0.97]$, it has at least four distinct colours after quantising each channel
to 8 levels, and the standard deviation of its opaque RGB is at least 18. These reject background patches, single-colour
blocks, and the near-empty crops that connected-component cutting produces at sheet margins.

\paragraph{Deduplication.}
Every surviving sprite is hashed by the exact bytes of its RGBA array together with its shape. Within the original
corpus this finds 4{,}888 extra copies. Because hold-out sets are defined by path, every copy of a held-out sprite is
also excluded from training, not only the copy that was drawn. Near-duplicates --- animation frames of the same
character --- are not removed from training: an ablation that collapsed them by an $8\times8$ luminance-plus-alpha
average hash lost more data than it gained diversity (training rows 64{,}990 to 49{,}984) and scored worse at equal
steps, so only the per-source cap below is applied.

\paragraph{Bucket policy.}
A sprite may feed a resolution bucket only if its native size is at most $1.5\times$ that bucket. This is the change
that stops 25--64\,px artwork from becoming a 16\,px target. It costs training rows: in the enlarged corpus, the
one whose training log we still have, it refuses 154{,}312 of 272{,}690 candidate (sprite, bucket) pairs. It is the
single largest contributor to the data-side gain.

\paragraph{Added material.}
14{,}999 sprites cut from three permissively licensed OpenGameArt collections (CC-BY-4.0, CC-BY-SA-3.0, OGA-BY-3.0)
pass the filters and cross-corpus deduplication. A second pass drops sheets whose member name matches a font, UI or
text pattern (367 sprites), collapses near-identical animation frames within a sheet, and caps any single source sheet
at 60 sprites so one animation cannot dominate a bucket. The survivors have a median palette of eight colours and
6{,}028 of them are natively 24\,px or smaller.

\paragraph{Captions.}
Sprites are captioned by a vision--language model shown numbered grids of 20 sprites on a checkerboard background,
which is an order of magnitude cheaper than one call per sprite. A text-only classifier then labels each caption
object / effect / fragment / glyph / junk; on the 2{,}980 sprites that also carry image-based labels the two agree
82.5\% of the time (90.5\% recall on objects), which is enough to use the labels for an ablation and not enough to
filter the training corpus with them.

\paragraph{Licences.}
The sprites added in this work carry per-file records of archive, member and licence. The older material predates that
bookkeeping; we release the manifests for what we added and the scripts that rebuild everything else from public
archives.

\section{Data interventions that did not earn a place}
\label{sup:data-abl}

Three attempts to fix the corpus at training time. None is in the final recipe; the second and third are reported
because of how they fail.

\paragraph{Content filtering.}
A text classifier labels every caption object / effect / fragment / glyph / junk. Dropping the 8{,}818 sprites it calls
fragment, glyph or junk (training rows 86{,}290 to 73{,}845) improves a 20k-step run from 165.6 to 156.7 native FD at
equal steps, a 5.4\% gain. That is real but small next to the sampler, it costs a full retraining, and the classifier
itself agrees with image-based labels only 82.5\% of the time, so we keep the unfiltered model.

\paragraph{Quality conditioning, and what it does to FD.}
A vision model rates each sprite's craft from 1 to 5; on the original corpus 25.2\% are rated 1--2 and 8.3\% are rated
5. Prepending the rating to the caption as a tag and asking for the top tag at sampling time is the standard
aesthetic-tag recipe and needs no architecture change. It works as intended in the only sense a person can check: the
samples are visibly more finished. Both distributional scores get worse --- 158.4 to 208.4 against the native
reference, and 186.5 to 227.5 against a reference restricted to the 719 sprites rated 4--5 --- because conditioning on
the best 8\% of the corpus narrows the output distribution, and FD measures distance to a distribution rather than to
its best part. We keep the main model unconditioned and report the knob, together with the disagreement: a reader who
looks at the samples and a reader who reads the table draw opposite conclusions, and both are reading their instrument
correctly.

\paragraph{More data.}
The body reports the headline of this one: two additions that look decisive at 20k steps (54.2 to 42.9 on the reported
configuration) and lose at 80k (39.3 $\pm$ 0.3 against 31.8 $\pm$ 0.4 over three seeds each). Table~\ref{tab:supdata}
gives both pilots and the full-length runs together. The two sources differ in what they add rather than how much: the
Kenney packs are flatter than anything else we have but only 41.8\% whole objects, while the CC-BY-3.0 bundle is
82.4\% whole objects. At 20k the second helps ten times as much as the first, which invited the reading that flatness
is what the decoder already supplies and shape is what it cannot. At 80k the effect is gone, and we do not offer a
mechanism: the natively-small fraction of the added material is 24.0\% against 28.6\% for what was already there,
nowhere near enough to explain a 24\% regression.

\begin{table}[h]
\caption{The enlarged corpus at two training lengths, 16\,px native FD. The 20k column ranks the corpora one way and
the 80k column the other.}
\label{tab:supdata}
\centering
\small
\begin{tabular}{@{}lcc@{}}
\toprule
Training corpus & 20k steps & 80k steps \\
\midrule
Base (86{,}290 rows)              & 54.2 & \textbf{31.8 $\pm$ 0.4} \\
{+ Kenney (98{,}590)}             & 53.2 & --- \\
{+ Kenney + CC-BY-3.0 (118{,}378)} & \textbf{42.9} & 39.3 $\pm$ 0.3 \\
\bottomrule
\end{tabular}
\end{table}

\section{Cross-resolution self-guidance: the negative result in full}
\label{sup:crossres}

A bucketed model offers a guidance reference with no extra weights: query the same network under a \emph{lower}
resolution label and extrapolate away from that prediction, $\tilde e = e_{\text{low}} + w\,(e_R - e_{\text{low}})$
with $e_{\text{low}} = \epsilon_\theta(x_t, t, c, b_{12})$. The reference is a prediction of what the same sprite would
look like drawn smaller, so the update pushes away from ``too coarse'' rather than away from ``unconditional''. We also
trained a one-pass version, which internalises the rule into the weights with a guidance-free objective so that a
single forward pass produces the guided prediction.

Under the conventional reference this was our strongest result. Under the native reference it is worth nothing, and
the internalised version is actively harmful (Table~\ref{tab:supcross}). The two columns disagree on every row and
they disagree in a consistent direction: the biased reference prefers more guidance, the native one prefers less.

\begin{table}[h]
\caption{16\,px, main model, 3{,}000 captions. Both references, same samples. The biased column orders the rows almost
exactly opposite to the native one.}
\label{tab:supcross}
\centering
\small
\begin{tabular}{@{}lcc@{}}
\toprule
Sampler & Biased $\downarrow$ & Native $\downarrow$ \\
\midrule
CFG 1.0 (no guidance)          & 19.0 & 113.2 \\
CFG 1.25                       & 15.2 & 113.2 \\
CFG 1.5                        & 13.6 & 117.9 \\
CFG 2                          & 12.2 & 131.1 \\
\midrule
cross-resolution, $w{=}1.5$    & 14.2 & 114.9 \\
composed reference, $w{=}1.5$  & 12.8 & 117.0 \\
internalised, one pass         & 10.4 & 170.3 \\
\midrule
ICG $w{=}1.25$                 & 15.5 & 94.2 \\
ICG $w{=}1.5$                  & 16.2 & 87.2 \\
ICG $w{=}2$                    & 20.7 & 87.1 \\
\bottomrule
\end{tabular}
\end{table}

The mechanism is visible in the per-image statistics. Guidance away from a lower-resolution prediction makes a sample
look like a sharpened version of a downscaled image, which is exactly what the biased reference is made of, so the
biased score improves monotonically with guidance weight. It does not reduce the colour count or increase flatness,
which is what separates the native reference from everything else, so the native score does not move. Independent-condition
guidance, which we did not invent, does help under both --- and it composes with palette projection, which is why the
final recipe uses it.

We report this at length because the two references would have led us to publish different papers. Had we scored only
against the conventional reference, the internalised one-pass rule --- 10.4, the best number in the table --- would
have been the contribution, and it is the single worst sampler in the column that measures distance to real pixel art.

\section{Ablations on the model rather than the decoder}
\label{sup:model-abl}

Before concluding that the remaining error is in the data and the decoder, we checked the three obvious properties of
the network. None of them explains it.

\paragraph{Text encoder.}
Swapping the frozen CLIP ViT-B/32 for ViT-L/14 (512 to 768 channels of cross-attention), same data, recipe and step
count: 164.0 against 165.6 native FD at 20k steps, a difference far inside what we can attribute, and caption
retrieval goes slightly \emph{down}, 20.6\% to 19.6\%. Whatever makes our samples draw the wrong object, it is not the
text tower, and a full retraining with the larger encoder is not worth 30 GPU-hours.

\paragraph{Training length.}
Continuing the main model past 80k steps flattens: 31.8 at 80k, 30.8 at 100k, 32.3 at 120k on the reported
configuration, all within a band of 1.5. We report 80k.

\paragraph{Capacity, and why we give the whole curve.}
A $2.2\times$ larger network (161.2M against 72.5M) leads early and loses late: 46.1 against 54.2 at 20k steps, 38.3
against 35.8 at 40k, and 38.5 against 31.8 at 80k. Read at 20k it is a clear win; read at 80k it is a clear loss. On
plain CFG the larger model is marginally better throughout (114.5 against 117.9), so the reversal is specific to the
configuration we report. The explanation we find plausible is that 2.2 times the parameters overfit a 50k-sprite corpus
sooner, but we did not test it, and the reason we show three points rather than one is that a single point would have
supported either conclusion --- the same trap that the corpus experiment in Section~\ref{sup:data-abl} fell into.

\section{Resolutions outside the target range}
\label{sup:res32}

The body covers 12--24\,px, the sizes the medium is actually used at. One rung above, at 32\,px, the ordering is the
same: CFG 1.5 scores 139.2 native FD, ICG $w{=}2$ 97.1, and ICG with an eight-colour projection 74.5 (twelve colours:
78.7). The best $k$ is 8, 4, 8, 6 and 8 at 12, 16, 20, 24 and 32\,px: it neither rises with the canvas nor tracks the real
median palette, and 16\,px --- the size we study most --- is the one that wants the fewest colours. Absolute values are not comparable across resolutions: each has
its own reference set and its own real-versus-real floor.

\section{Alternative reference sets}
\label{sup:refs}

\paragraph{Object-only reference.}
About a fifth of the corpus is not a whole object: effects, fragments of larger drawings, glyphs. Restricting the
reference to the 825 held-out sprites a classifier calls objects, and the prompts to the matching 2{,}181, gives 37.1
for the reported configuration. The ordering of methods is unchanged; the absolute value moves because the reference
population did.

\paragraph{Clean subset.}
Removing the 41.4\% of held-out prompts that have a near-duplicate twin in training leaves 1{,}759 prompts and moves
the reported configuration from 31.8 to 29.0, with every other row moving by at most 3.0 and no ordering changing.

\paragraph{A second feature space.}
Inception-v3 clean-FID and KID on the same samples and the same native reference preserve the ordering, controls
included: 23.79 for the projection, 23.28 at $k{=}6$, 26.27 without it, 29.37 for CFG, 36.05 for one global palette,
66.63 for the median filter, 10.54 for the real held-out sprites. The separation shrinks sharply --- a $2.8\times$ gap
between projecting and not projecting under DINOv2 is $1.10\times$ here, and $k{=}4$ and $k{=}6$ swap --- which is what
an encoder whose input is resized to 299\,px would do to 16\,px hard edges. We read it as a check on the ranking and
not as confirmation of the effect size.

\section{Prompt alignment in full}
\label{sup:clip}

CLIP ViT-B/32 image--text similarity ($100\cos$) and retrieval of the own caption among 100 distractors, 3{,}000
held-out captions, 16\,px unless stated. ICG $w{=}2$ alone: 27.60 / 22.8\%. Projection at $k{=}6$: 27.38 / 22.0\%. At
$k{=}4$: 27.07 / 19.2\%. The real held-out sprites the captions were written from: 27.66 / 22.6\%, so at $k{=}6$ our
samples are as retrievable from their captions as real sprites are. At 20\,px the projection costs 1.9 points (26.7\%
to 24.8\%) and at 24\,px 1.9 points (27.6\% to 25.7\%).

Two cautions about this measurement. The distractor draw is re-seeded per directory: with a single shared stream, a
set's retrieval score depended on which other sets were scored in the same invocation, which cost us about a point in
an earlier version of these numbers. And a real-sprite baseline is only meaningful if the sprites are paired with the
captions that describe them --- an unpaired reference scores 19.68 / 1.1\%, chance by construction, and reading that
as evidence that real pixel art is hard to retrieve would be a mistake.

\section{Licences and release}
\label{sup:licence}

The sprites added in this work carry per-file records of archive, member file and licence (CC-BY-4.0, CC-BY-SA-3.0,
OGA-BY-3.0 for the OpenGameArt material; CC0-1.0 for the Kenney packs used in the data ablation). The older material
predates that bookkeeping; we release the manifests for what we added, together with the scripts that rebuild the rest
from public archives and a table mapping every number in the paper to the script that produces it.

\end{document}
